\documentclass[12pt]{article}
\usepackage{amsmath}
\usepackage{amssymb}
\usepackage{geometry}
\usepackage{array}
\usepackage{booktabs}

\geometry{margin=1in}

\title{\textbf{Rational Exponents}}
\author{}
\date{}

\begin{document}

\maketitle
\thispagestyle{empty}

% ─────────────────────────────────────────
\section*{Page 2 — Building Intuition}

\textbf{Recall:}
\[
2 + 2 + 2 + 2 + 2 = 5 \times 2
\]
\[
\Rightarrow \quad \text{Multiplication is \underline{repeated addition}}
\]

\bigskip

\[
3 \cdot 3 \cdot 3 \cdot 3 \cdot 3 \cdot 3 = 3^{6}
\]
\[
\Rightarrow \quad \text{Exponent is \underline{repeated multiplication}}
\]

\bigskip

\textbf{My Favourite Exponent:}
\[
\underbrace{2^{10}}_{\text{base}} = \underbrace{2 \cdot 2 \cdots 2}_{10} = 1024 \quad \leftarrow \text{product}
\]

\bigskip

\begin{center}
\begin{tabular}{lll}
\toprule
\textbf{Type} & \textbf{Goal} & \textbf{Example} \\
\midrule
Type 1 (learned) & Calculate the \textit{product} & $2^{10} = \underline{\quad\quad}$ \\[6pt]
Type 2 (today)   & Find the \textit{base}          & $\boxed{\phantom{0}}^{10} = 1024$ \\[6pt]
Type 3 (future)  & Find the \textit{power/exponent} & $10 = \log_{2} 1024$ \\
\bottomrule
\end{tabular}
\end{center}

% ─────────────────────────────────────────
\newpage
\section*{Page 3 — Finding the Base}

\textbf{Analogy with multiplication:}
\[
2 \times 3 = 6
\quad \Rightarrow \quad
\square \cdot 3 = 6
\quad \Rightarrow \quad
\square = 6 \div 3 \quad \text{(long division)}
\]

\bigskip

\textbf{Same idea for exponents:}
\[
2^{10} = 1024
\quad \Rightarrow \quad
\square^{10} = 1024
\]
\[
\underbrace{x \cdot x \cdot x \cdot x \cdots x}_{10} = 1024
\]

\bigskip

\textbf{Answer:}
\[
\boxed{2 = \sqrt[10]{1024}}
\]

% ─────────────────────────────────────────
\newpage
\section*{Page 4 — Radical Notation}

We introduce the \textbf{radical sign}:
\[
\sqrt[n]{\,a\,}
\]

\begin{center}
\begin{tabular}{ll}
$n$ & : \textbf{index} \\[4pt]
$a$ & : \textbf{radicand} \\
\end{tabular}
\end{center}

% ─────────────────────────────────────────
\newpage
\section*{Page 5 — Inverse Relationship}

\textbf{Recall from multiplication:}
\[
3 \cdot 4 = 12
\quad \Rightarrow \quad
4 = 12 \div 3
\quad \Rightarrow \quad
3 = 12 \div 4
\]

\bigskip

\textbf{Same idea for exponents:}
\[
2^{2} = 4 \quad \Rightarrow \quad 2 = \sqrt{4}
\]
\[
3^{2} = 9 \quad \Rightarrow \quad 3 = \sqrt{9}
\]
\[
2^{3} = 8 \quad \Rightarrow \quad 2 = \sqrt[3]{8}
\]
\[
3^{3} = 27 \quad \Rightarrow \quad 3 = \sqrt[3]{27}
\]

\bigskip

\textbf{Convention:}
\[
\sqrt{\phantom{0}} = \sqrt[2]{\phantom{0}}
\quad \text{(index 2 is invisible — most commonly used root)}
\]
\[
2 = 2^{1}
\quad \text{(similarly, exponent 1 is invisible)}
\]

% ─────────────────────────────────────────
\newpage
\section*{Page 6 — Fill In}

\[
2^{5} = \boxed{32}
\quad \Rightarrow \quad
2 = \sqrt[5]{32}
\]

\[
3^{4} = \boxed{81}
\quad \Rightarrow \quad
3 = \sqrt[4]{81}
\]

\[
4^{2} = \boxed{\phantom{000}}
\quad \Rightarrow \quad
4 = \sqrt{\phantom{000}}
\]

\[
2^{10} = \boxed{\phantom{000}}
\quad \Rightarrow \quad
2 = \sqrt[10]{\phantom{000}}
\]

\end{document}
